% Auto-generated by scripts/f5_vl2_plm_amplitude_decay.py
% Do not edit manually.

% PLM central-slope + upwind semi-discrete operator
\begin{equation}\label{eq:F5_Lhat}
\hat{L}(\xi) = -\frac{a}{h}\bigl(1 + \tfrac{i}{2}\sin\xi\bigr)\,\bigl(1 - e^{-i\xi}\bigr),\quad \xi = k h
\end{equation}

% Midpoint-RK2 amplification factor
\begin{equation}\label{eq:F5_g}
g(\xi; \nu) = 1 + \nu\hat{L} + (\nu\hat{L})^2/2,\quad \nu = a\Delta t/h
\end{equation}

% 2nd-order signature: leading dissipation is O(ξ⁴)
\begin{equation}\label{eq:F5_amp}
|g(\xi;\nu)|^2 = 1 + c_4(\nu)\,\xi^4 + \mathcal{O}(\xi^6),\quad c_4(\nu) < 0\ \text{for }|\nu| < 1
\end{equation}

% Amplitude retention over fixed time
\begin{equation}\label{eq:F5_retention}
\frac{\mathrm{amp}(t)}{\mathrm{amp}(0)} = |g|^{N_{\text{step}}},\quad N_{\text{step}} = \frac{t\,a}{\nu\,h} \propto h^{-1}
\end{equation}

% Decay rate scales as h³ (NOT h²)
\begin{equation}\label{eq:F5_h3}
\gamma_{\text{num}} = -\frac{\ln|g|^{N_{\text{step}}}}{t}\ \propto\ h^3 \ \Longleftrightarrow\ N^{-3}
\end{equation}

% Two-resolution inversion for 2nd-order schemes (CORRECTED F4)
\begin{equation}\label{eq:F5_p3}
p \equiv \frac{\log(\gamma_1/\gamma_2)}{\log(N_2/N_1)} = 3\quad \text{(for amplitude-retention measurement)}
\end{equation}

% Comparison: L¹ error convergence vs amplitude retention
\begin{equation}\label{eq:F5_L1_vs_gamma}
\text{L}^1\text{-error}(\mathrm{IC}) \propto h^2\quad \text{(standard 2nd-order convergence)}\\ \gamma_{\text{num}}(\mathrm{amp}) \propto h^3\quad \text{(amplitude-retention rate; this benchmark)}
\end{equation}

% Why F4 was wrong
\begin{equation}\label{eq:F5_F4_fix}
F4\text{-order stated } p = 2 \text{ based on } \nu_{\text{eff}}\propto h^2\text{ analysis.  But } \gamma = \nu_{\text{eff}} k^2/2 \propto h^2\text{ holds only per step.}\\ \text{Total decay over fixed }t\text{ picks up extra }1/h\text{ from }N_{\text{step}}.\\ \text{Corrected expectation: } p = 3.
\end{equation}
